
\documentclass[a4paper,10pt]{article}
\usepackage{amsmath}
\usepackage{ctex} % 中文支持
\usepackage{geometry}
\usepackage{float}

% 设置页边距
\geometry{top=1in, bottom=1in, left=1in, right=1in}

% 页眉设置
\usepackage{fancyhdr}
\pagestyle{fancy}
\fancyhead[L]{实验预习报告} % 左侧页眉内容
\fancyhead[C]{} % 中间页眉内容
\fancyhead[R]{韩初晓} % 右侧页眉内容，可以填写你的名字
\fancyfoot[C]{\thepage} % 页脚中间显示页码

% 正文开始
\begin{document}

\title{气轨实验预习报告} % 标题
\author{韩初晓} % 作者名
\date{\today} % 日期，使用当前日期
\maketitle % 生成标题

% 下面是正文内容
\newpage


\section{简谐振动的原理}

在气垫导轨上，滑块与两弹簧连接，滑块做简谐振动。假设滑块质量为 $m_1$，处于平衡位置时，弹簧的伸长量为 $x_0$。当滑块位移为 $x$ 时，弹簧分别施加的力为 $-k_1(x + x_0)$ 和 $-k_1(x - x_0)$，其中 $k_1$ 是弹簧的劲度系数。由牛顿第二定律，得运动方程：

\[
m_1 \ddot{x} = -k(x)
\]

其中 $k = 2k_1$。此方程的解为：

\[
x(t) = A \sin(\omega_0 t + \varphi_0)
\]

其中 $A$ 为振幅，$\omega_0 = \sqrt{k/m}$ 为固有频率。

\subsection{振动周期与质量关系}

振动系统的有效质量为 $m = m_0 + m_1$，其中 $m_0$ 是弹簧的有效质量，$m_1$ 为滑块质量。振动周期 $T$ 与 $\omega_0$ 的关系为：

\[
T = 2\pi \sqrt{\frac{m}{k}} = 2\pi \sqrt{\frac{m_1 + m_0}{k}}
\]

平方后得：

\[
T^2 = \frac{4\pi^2 (m_1 + m_0)}{k}
\]

通过改变 $m_1$，测得不同 $T$ 值，并绘制 $T^2$ 与 $m_1$ 的关系曲线，斜率为 $\frac{4\pi^2}{k}$，可以计算出 $k$ 的值，截距得到 $m_0$。

\section{运动学特征与能量分析}

\subsection{速度与位置关系}

对振动方程求导，得到瞬时速度：

\[
v(t) = A \omega_0 \cos(\omega_0 t + \varphi_0)
\]

速度与位置满足关系：

\[
v^2 = \omega_0^2 (A^2 - x^2)
\]

当 $x = A$ 时，$v = 0$；当 $x = 0$ 时，$v = \pm A \omega_0$，此时速度最大。

\subsection{机械能守恒}

在任何时刻，系统的动能和势能分别为：

\[
E_k = \frac{1}{2} (m_1 + m_0) v^2, \quad E_p = \frac{1}{2} k x^2
\]

系统的总机械能 $E$ 为：

\[
E = E_k + E_p = \frac{1}{2} k A^2
\]

由于总能量守恒，系统的机械能在振动过程中保持不变。

\section{瞬时速度的测定}

瞬时速度是物体在某一时刻的速度，理论定义为：

\[
v_{\text{瞬}} = \lim_{\Delta t \to 0} \frac{\Delta s}{\Delta t}
\]

实验中，通过测量小时间间隔内的平均速度来间接测定瞬时速度。假设滑块从 A 点下滑，经过距离 $\Delta s$，时间为 $\Delta t$，其平均速度为：

\[
v = \frac{\Delta s}{\Delta t}
\]

当 $\Delta t \to 0$ 时，平均速度的极限值即为瞬时速度 $v_0$。实验中通过改变挡光距离 $\Delta s$，测量不同时间间隔内的平均速度，并绘制 $v$ 与 $\Delta t$ 的图线，外推至 $\Delta t = 0$ 处，得到瞬时速度。


\section{实验操作}

1. 学会使用光电门测量滑块的速度，并使用它来测量振动周期。

2. 调节气垫导轨至水平状态，测量滑块在任意两点的速度变化，验证气垫导轨是否水平。

3. 分别取振幅 $A = 10.0, 20.0, 30.0, 40.0$ cm，测量相应的振动周期，分析周期与振幅的关系。

4. 在滑块上加骑码（铁片），固定振幅 $A = 40.0$ cm，增加骑码测量周期 $T$。绘制 $T^2$ 对质量 $m$ 的图，求得弹簧劲度系数 $k$ 和有效质量 $m_0$。

5. 安装 U 型挡光片，测量速度，绘制 $v^2$ 对 $x^2$ 的图，进行拟合，求得斜率和截距，验证 $- \omega_0^2$ 与 $A^2 \omega_0^2$ 是否为斜率和截距。

6. 固定振幅 $A = 40.0$ cm，测量不同位置 $x$ 处的速度，计算动能与势能，并比较各位置的机械能，验证机械能是否守恒。

7. 改变挡光片距离，测量滑块从静止下滑时的平均速度，绘制 $v$ 对 $\Delta t$ 和 $\Delta s$ 的图，使用外推法求瞬时速度。

8. 改变气轨的倾斜角度 $\theta$，重复上述实验。

9. 改变 $A$ 点到 $P$ 点的距离 $l$（设置为 $60$ cm），重复实验。

\section{注意事项}

\begin{enumerate}
  \item 先开气源前后放滑块，实验结束时先取滑块再关气源。
  \item 滑块需轻拿轻放，挂弹簧时手扶住滑块。
  \item 测量周期时使用条形挡光片，测量速度时使用 U 形挡光片。
  \item 调节挡光片或放骑码时，必须将滑块取下后再操作。
  \item 测定瞬时速度时，前后安装 U 形挡光片确保配重均衡。
  \item 测定瞬时速度时，滑块应静止释放，滑过光电门后及时扶住滑块，避免撞到弹簧。
\end{enumerate}


\end{document}
